\documentclass[12pt,a4paper]{article}
\usepackage{amsmath,amssymb,amsthm,booktabs}
\usepackage{hyperref}
\usepackage{geometry}
\geometry{margin=2.5cm}
\usepackage{enumitem}

\newtheorem{theorem}{Theorem}[section]
\newtheorem{lemma}[theorem]{Lemma}
\newtheorem{corollary}[theorem]{Corollary}
\newtheorem{proposition}[theorem]{Proposition}
\theoremstyle{definition}
\newtheorem{definition}[theorem]{Definition}
\theoremstyle{remark}
\newtheorem{remark}[theorem]{Remark}

\title{The EDGES 21\,cm Anomaly:\\
Status, Analysis, and RS Perspective}

\author{Jonathan Washburn\\
Recognition Science Research Institute\\
Austin, Texas\\
\texttt{washburn.jonathan@gmail.com}}

\date{March 2026}

\begin{document}
\maketitle

\begin{abstract}
We analyze the EDGES 21\,cm absorption anomaly within Recognition
Science (RS).  EDGES (2018) reported a global 21\,cm absorption
trough at $z \approx 17$ with amplitude $T_{21} \approx
-500^{+200}_{-500}~\mathrm{mK}$, roughly $2$--$3\times$ deeper than
the standard $\Lambda$CDM prediction of $T_{21} \approx
-200~\mathrm{mK}$~\cite{EDGES}.  However, the SARAS-3 experiment
(2022) reported no detection of a cosmological 21\,cm signal at the
EDGES amplitude, and determined that the EDGES spectral profile is
inconsistent with a cosmological origin at $95.3\%$ confidence,
suggesting an instrumental or foreground systematic in the EDGES
measurement~\cite{SARAS3}.  The experimental status of the anomaly
is therefore \emph{ambiguous}: the EDGES excess may not be a real
cosmological signal.  We present the RS analysis under both scenarios.
If the signal is real, RS offers an enhanced cooling mechanism via
the 8-tick recognition cycle.  If it is instrumental, the RS
prediction is $T_{21} \approx -200~\mathrm{mK}$, consistent with
standard adiabatic cooling.  No new particles are required under
either interpretation.
\end{abstract}

%% ============================================================
\section{Recognition Science Framework}
\label{sec:framework}
%% ============================================================

The RS substrate-cooling interpretation of the EDGES signal rests on
the unique cost functional determined by the following axioms.

\begin{definition}[RS Cost Functional]
For $x > 0$: $J(x) = \tfrac{1}{2}(x + x^{-1}) - 1$.
\end{definition}

\noindent\textbf{Axiom A1 (Normalization):} $J(1) = 0$.\\
\textbf{Axiom A2 (Recognition Composition Law, RCL):}
$J(xy)+J(x/y)=2J(x)J(y)+2J(x)+2J(y)$ for all $x,y>0$.\\
\textbf{Axiom A3 (Calibration):} $J''_{\log}(0) = 1$, where
$J_{\log}(t) := J(e^t)$.

\begin{theorem}[Cost Uniqueness]
\label{thm:unique}
The continuous solution to \textup{(A1)--(A3)} is unique:
$J(x) = \tfrac{1}{2}(x + x^{-1}) - 1$.
\end{theorem}

\begin{proof}[Proof sketch]
Set $H(t) = J_{\log}(t)+1$.  Axiom A2 transforms into the d'Alembert
equation $H(s+t)+H(s-t)=2H(s)H(t)$.  By the Acz\'el classification
theorem~\cite{Aczel1966}, the unique continuous solution with $H(0)=1$
and $H''(0)=1$ is $H(t)=\cosh t$, giving
$J(x)=\tfrac{1}{2}(x+x^{-1})-1$.
\end{proof}

\noindent The recognition ledger executes one complete update cycle every
8 ticks ($2^D = 8$, with $D = 3$ forced by $2^D = 8$).  At each tick,
the ledger minimizes the total $J$-cost.  In a cooling context, this
provides a channel for kinetic energy to transfer from the baryon gas
into substrate coherence.  The coupling efficiency $\epsilon_{\rm sub}$
is the RS parameter governing this transfer; its first-principles
derivation is an open problem identified in Section~\ref{sec:analysis}.

%% ============================================================
\section{Introduction}
\label{sec:intro}
%% ============================================================

Recognition Science~\cite{RS_Axioms} (RS) provides a substrate cooling
mechanism via the 8-tick recognition cycle that may enhance the 21\,cm
absorption if the EDGES signal is confirmed cosmological.

The 21\,cm hyperfine transition of neutral hydrogen provides a probe
of the intergalactic medium (IGM) during cosmic dawn ($z \sim 15$--$25$).
The brightness temperature contrast against the CMB is:
\begin{equation}
  T_{21} \propto \left(1 - \frac{T_{\mathrm{CMB}}}{T_{\mathrm{spin}}}\right),
\end{equation}
where $T_{\mathrm{spin}}$ is the hydrogen spin temperature.  If the
gas is colder than the CMB ($T_{\mathrm{spin}} < T_{\mathrm{CMB}}$),
the signal is in absorption.  Standard adiabatic cooling predicts
$T_{21} \approx -200~\mathrm{mK}$ at $z \approx 17$.

\subsection{EDGES Result (2018)}

EDGES measured the global 21\,cm signal at $78~\mathrm{MHz}$
($z \approx 17$) and found $T_{21} = -500^{+200}_{-500}~\mathrm{mK}$,
roughly $2$--$3\times$ deeper than the standard prediction~\cite{EDGES}.
If genuine, this would require either the gas to be colder (beyond
standard adiabatic cooling) or the background radiation to be hotter.

\subsection{SARAS-3 Result (2022)}

The SARAS-3 experiment independently measured the global 21\,cm signal
using a different antenna design and site.  Reporting in 2022,
SARAS-3 found~\cite{SARAS3}:
\begin{itemize}[nosep]
  \item No detection of a cosmological signal at the EDGES amplitude.
  \item A Bayesian analysis disfavoring the EDGES spectral profile as
  a cosmological signal at $95.3\%$ confidence.
  \item Consistency with the SARAS-3 data being dominated by
  foreground contamination and/or instrumental effects.
\end{itemize}

This result has cast significant doubt on whether the EDGES signal
has a cosmological origin.  The experimental status of the 21\,cm
anomaly at $z \approx 17$ is currently unresolved.

%% ============================================================
\section{RS Analysis}
\label{sec:analysis}
%% ============================================================

\subsection{Scenario A: EDGES Signal is Real}

If the EDGES absorption trough is confirmed by future experiments, RS
offers a substrate cooling mechanism.

\begin{theorem}[Standard Adiabatic Cooling Prediction]
\label{thm:adiabatic}
In the standard cosmological model, gas temperature evolves as
$T_{\mathrm{gas}} \propto (1+z)^2$ after thermal decoupling from the
CMB at $z \sim 150$.  At $z \approx 17$:
\begin{equation}
  T_{\mathrm{gas}}^{\mathrm{std}} \approx T_{\mathrm{CMB}}^{z=150}
  \times \left(\frac{1+17}{1+150}\right)^2 \approx 5.5~\mathrm{K}.
\end{equation}
The CMB temperature at $z = 17$ is $T_{\mathrm{CMB}} = 2.725 \times 18
= 49.1~\mathrm{K}$, giving an absorption depth $T_{21} \approx
-200~\mathrm{mK}$.
\end{theorem}

\begin{proposition}[Substrate-Enhanced Cooling]
\label{prop:cooling}
In RS, the 8-tick recognition cycle provides an additional energy
extraction channel: at each tick, the ledger minimizes the total
$J$-cost, transferring kinetic energy from the gas into substrate
coherence.  This substrate coupling would produce additional cooling
beyond adiabatic:
\begin{equation}
  T_{\mathrm{gas}}^{\mathrm{RS}} \approx \frac{T_{\mathrm{gas}}^{\mathrm{std}}}
  {1 + \epsilon_{\mathrm{sub}}},
\end{equation}
where $\epsilon_{\mathrm{sub}}$ is the substrate coupling efficiency.
For $\epsilon_{\mathrm{sub}} \sim 1$--$2$, this gives an additional
factor of $2$--$3$ cooling, producing $T_{21} \approx -400$ to
$-600~\mathrm{mK}$.
\end{proposition}

\begin{remark}
The coupling efficiency $\epsilon_{\mathrm{sub}}$ is not derived from
first principles in the present work.  The RS framework predicts that
substrate cooling exists, but the quantitative rate depends on the
coupling between the gas and the ledger at $z \sim 17$, which
requires a detailed calculation.  The factor of $2$--$3$ cooling is
stated as an order-of-magnitude target, not a parameter-free prediction.
\end{remark}

\subsection{Scenario B: EDGES Signal is Instrumental}

If SARAS-3 and future experiments confirm that the EDGES excess is
a systematic artifact, the RS prediction is consistent with the
standard adiabatic result:

\begin{remark}[RS Default Prediction]
\label{rem:standard}
If no enhanced cooling mechanism is activated at $z \sim 17$
(i.e., $\epsilon_{\mathrm{sub}} \ll 1$ at this epoch), the RS
prediction reduces to the standard adiabatic result
$T_{21} \approx -200~\mathrm{mK}$.  The enhanced substrate cooling
is an optional feature of the RS framework, not a mandatory
prediction.  If SARAS-3 and future experiments confirm that the
EDGES excess is an instrumental artifact, this default applies.
\end{remark}

%% ============================================================
\section{Experimental Comparison}
\label{sec:compare}
%% ============================================================

\begin{center}
\renewcommand{\arraystretch}{1.3}
\begin{tabular}{lccc}
\toprule
\textbf{Observable} & \textbf{EDGES (2018)} & \textbf{SARAS-3 (2022)} & \textbf{RS}\\
\midrule
$T_{21}$ at $z \approx 17$ & $-500^{+200}_{-500}$ mK & Not detected & $-200$ to $-600$ mK\\
Signal cosmological? & Claimed & Disfavored & Either scenario\\
No new particles & No DM needed & --- & Confirmed\\
\bottomrule
\end{tabular}
\end{center}

%% ============================================================
\section{Predictions and Falsifiers}
\label{sec:predict}
%% ============================================================

\begin{enumerate}
  \item \textbf{Key experiment: REACH and PRIZM.}  The REACH
  (Radio Experiment for the Analysis of Cosmic Hydrogen) and PRIZM
  experiments are designed to independently measure the global 21\,cm
  signal with different instrument designs.  Their results will be
  decisive.  If they confirm the EDGES signal at the same amplitude
  and phase, the anomaly is likely cosmological.  If they find no
  signal, the SARAS-3 interpretation is correct.

  \item \textbf{Millicharged dark matter exclusion.}  Whether or not
  the EDGES signal is cosmological, RS predicts no millicharged dark
  matter particles.  Dedicated searches (SENSEI, MilliQan) will yield
  null results.

  \item \textbf{$\varphi$-structure in confirmed signal.}  If
  future experiments confirm a deep 21\,cm trough, the RS substrate
  cooling mechanism predicts that the absorption profile will show
  subtle $\varphi$-ladder structure (modulation of the trough shape
  at frequency ratios $\nu_n/\nu_0 = \varphi^n$), distinguishable
  from smooth adiabatic profiles.

  \item \textbf{Falsifier for RS enhanced cooling.}  If confirmed
  experiments establish $T_{21} = -200 \pm 30~\mathrm{mK}$ at
  $z \approx 17$ (standard prediction), the substrate-enhanced cooling
  hypothesis is falsified.
\end{enumerate}

%% ============================================================
\section{Conclusion}
\label{sec:conclusion}
%% ============================================================

The EDGES 21\,cm anomaly is experimentally unresolved.  SARAS-3
(2022) has cast significant doubt on whether the EDGES signal has
a cosmological origin.  RS accommodates both outcomes: if the signal
is real, substrate-enhanced cooling via the 8-tick cycle provides a
qualitative mechanism; if instrumental, the standard adiabatic
prediction ($T_{21} \approx -200~\mathrm{mK}$) applies.  No new
particles (millicharged dark matter, primordial magnetic fields) are
required under either interpretation.  The REACH and PRIZM experiments
will provide the decisive test.

%% ============================================================
\begin{thebibliography}{99}

\bibitem{EDGES}
J.~D.~Bowman et al.,
``An absorption profile centred at 78 megahertz in the
sky-averaged spectrum,''
Nature \textbf{555}, 67 (2018).

\bibitem{SARAS3}
N.~Rao et al.\ (SARAS-3 Collaboration),
``Absorption Feature in the Sky-Averaged Radio Spectrum between 55 and
85\,MHz: Constraints from the SARAS\,3 Radiometer,''
Nature Astronomy \textbf{6}, 1473 (2022).

\bibitem{Barkana2018}
R.~Barkana,
``Possible interaction between baryons and dark-matter particles
revealed by the first stars,''
Nature \textbf{555}, 71 (2018).

\bibitem{RS_Axioms}
J.~Washburn, ``The Algebra of Reality: A Recognition Science Derivation
of Physical Law,'' \emph{Axioms} \textbf{15}(2), 90 (2026).
\texttt{doi:10.3390/axioms15020090}

\bibitem{Aczel1966}
J.~Acz\'el, \emph{Lectures on Functional Equations and Their
Applications} (Academic Press, 1966).

\end{thebibliography}

\end{document}
